% This is samplepaper.tex, a sample chapter demonstrating the
% LLNCS macro package for Springer Computer Science proceedings;
% Version 2.20 of 2017/10/04
%
\documentclass[runningheads]{llncs}
%
\usepackage{graphicx}
% Used for displaying a sample figure. If possible, figure files should
% be included in EPS format.
%
% If you use the hyperref package, please uncomment the following line
% to display URLs in blue roman font according to Springer's eBook style:
% \renewcommand\UrlFont{\color{blue}\rmfamily}

\usepackage{longtable}
\usepackage{booktabs}
\usepackage{array}
% Used for table definitions in the "Data Collection" section.

\usepackage{amsmath}
% Used for mathematical equations in the "Data Cleaning" section.

\begin{document}
%
\title{Anaylizing Historical Stock Data using Machine Learning for Dummies}
%
%\titlerunning{Abbreviated paper title}
% If the paper title is too long for the running head, you can set
% an abbreviated paper title here
%
\author{Caleb R. Sellinger}
%
\authorrunning{C. Sellinger}
% First names are abbreviated in the running head.
% If there are more than two authors, 'et al.' is used.
%
\institute{Northwest Missouri State University, Maryville MO 64468, USA \\
\email{S577274@nwmissouri.edu}}
%
\maketitle              % typeset the header of the contribution
%
\begin{abstract}
The abstract should briefly summarize the contents of the paper in
150--250 words. The abstract should briefly summarize the contents of the paper in
150--250 words. The abstract should briefly summarize the contents of the paper in
150--250 words. The abstract should briefly summarize the contents of the paper in
150--250 words.

\keywords{stock market \and machine learning \and data analysis \and predictive modeling}
\end{abstract}
%
%
%
\section{Introduction}

I am going to analyze historical stock data in order to more easily make decisions within
the stock market for personal investment accounts.

I plan to source data from YahooFinance via a web scraper, extracting and storing it in a
file to use for analysis.\cite{YahooFinance}

The stock market is a very complex system. There is a lot of business terminology and acumen
that makes it difficult for the average person to make sound financial decisions that effectively
invest in there future. Therefore, my goal is to provide create a project that makes it more
approachable. Possible ideas are giving each stock a score or returning top X stocks based on the
strategy the individual wants use (dividend investment, stock growth, tax-loss harvesting, etc.).
These results will help the user to make a clear decision on what to invest in for their personal
accounts on apps like Robinhood, Fidelity, Webull, etc. Side goal, this may also provide a free
and open-source platform for stock market discovery or learning, no need to pay for tokens or
other AI agents.

\section{Project Methodology}

Process:
\begin{enumerate}

\item Web scraper extracts data from historical data from YahooFinance, including historical prices, dividend payouts, stock splits, quarterly income statements, quarterly balance sheets, cash flow, and quarterly statistics such as market cap, p/e ratios, etc. This data is stored in a csv file for analysis.
\item Analyze stock data in csv format. Gather key statistics that appear viable and trend historical data to identify growth potential. Possible ranking/scoring system for each data point.
\item Evaluate stock data based on user strategy. Predict viable outcomes of future stock performance against historical data and key statistics. Rank/score viable stocks for investment opportunity.
\item Process results in a digestible fashion in order for the user to make clear decisions.

\end{enumerate}

\section{Approach and Scope}

Using Jupyter Notebook and Python, I will create a web scraper to extract historical stock data from YahooFinance. The data will be stored in a csv file, and then analyzed using some sort of machine learning algorithm, further data investigation needed, to identify trends and patterns that can help predict future stock performance. The scope of this project is limited to analyzing historical stock data and providing insights for personal investment decisions. The project will focus on providing a user-friendly output for individuals to access and interpret the data, with the goal of making stock market analysis more approachable for those without extensive financial knowledge. Minimum viable product will be an output score of single stock analysis. Future work may include analysis of all stocks in the S\&P 500 or other indexes to produce a list of top X stocks based on user strategy. Predictions may vary in accuracy and should not be taken as financial advice and should be used for educational purposes only. The project will not provide any financial advice or recommendations, and users should always consult with a financial advisor before making any investment decisions.

\section{Data Collection}

The data is sourced from YahooFinance\cite{YahooFinance} using the requests and lxml library in Python to send a GET request, extract the DOM and parse the XML, using XPath, to scrape the data from their HTML elements. The base URL and headers are generated by a custom function to get to the main page of YahooFinance, and then scrape functions for each file are used to create an HTTP GET request for the specific page containing the data, store the response in a string variable, scrape XML using XPath method to find the <div> or <table> tag that the data is in, parse the element's child tags, and store data in a DataFrame. Then, the data is appropriately cleaned and stored in csv files.

For the historical dividend data, the data is scraped from the historical data tab, filtered by dividend only and date range set to max. For the unpaid version of YahooFinance, this date range is limited to 20 years of historical data. This data only has two attributes, the date and dividend amount that was paid on that date. The number of records varies by stock, as some stocks may pay dividends monthly, quarterly, or annually, or IPO'd less than 20 years ago.

For the key statistics data, the data is scraped from the statistics tab, which has quarterly summaries dating back five quarters and the current summary as dates. The number of records is nine. The program transposes the data to have the attributes as columns and the dates as rows. The values in each row, that are not dates, are converted from a string to a float, and the dates are converted to ad datetime object to allow for proper analysis.

For the income statement and cash flow data, the data is scraped from the financials tab and their respective tables. The data only includes quarterly summaries dating back four quarters and the current summary. The number of records for the income statement is 32 and the cash flow is 10.

The goal of collecting data in this format, is to provide a comprehensive view of a stock's financial performance over time without having to manually collect and organize the data. YahooFinance keeps up-to-date, live data from their servicer. Not all attributes will be used the machine learning algorithm, but the user may be interested in the available information. There are several limitations in the data, mainly, the data is limited to what the unpaid version of YahooFinance provides, which is limited to 20 years of historical data for dividends, five quarters of key statistics, and four quarters of income statement and cash flow data. Additionally, the data may not be complete or accurate, as it is sourced from a third-party website and may be subject to changes or omissions. Noted on the financials tab, YahooFinance is ``moving the Financials tab to a new, more detailed dataset from Morningstar\ldots''\cite{YahooFinanceHelp} on November 9\textsuperscript{th}, 2026.

\section{Data Cleaning}

\section{Analysis and Evaluation}

% Net Income Common Stockholders / (Total Assets - Total Liabilities) *Pulled from Balance Sheet = RoE *r on DCF model
% Linear Regression model, scikit-learn, use dates for x and free cash flow for y

% Net Income Common Stockholders is net income - preferred dividends. I am using this over the raw income, in order to get a more accurate representation of the company's profitability for common shareholders. This is important because preferred shareholders have a higher claim on the company's assets and earnings, and their dividends are typically fixed, while common shareholders have a residual claim on the company's earnings and may receive variable dividends based on the company's performance. By using Net Income Common Stockholders, we can better understand how much profit is available to common shareholders after accounting for the fixed obligations to preferred shareholders. (tldr; treating preferred dividends as an expense)

\section{Modeling}

The data curation process begins by extracting financial data from YahooFinance and storing the information in individual dataframes. Once the data has been extracted, a separate cleaning function is applied to each dataframe to ensure consistency, accuracy, and usability. These functions include \texttt{clean\_}\allowbreak\texttt{dividends}, \texttt{clean\_}\allowbreak\texttt{statistics}, \texttt{clean\_}\allowbreak\texttt{financials}, \texttt{clean\_}\allowbreak\texttt{cashflow}, and \texttt{clean\_}\allowbreak\texttt{balance}\allowbreak\texttt{sheet}. Each function performs several cleaning operations, such as renaming or removing unnecessary columns, replacing null or blank values, removing or modifying unwanted characters within data values, such as commas, and converting data into appropriate data types. This conversion is particularly important because much of the information retrieved from YahooFinance is initially stored in string format, which can make numerical analysis difficult. Additionally, the statistics and cash flow dataframes are transposed to improve their structure and make financial information across different time periods easier to interpret and analyze. The methods used for cleaning the data are built into the pandas and datetime libraries. Pandas methods include \texttt{.replace}, \texttt{.rename}, \texttt{.T} (for transposing a table), \texttt{.drop}, and \texttt{.reset\_index}. Datetime methods include \texttt{.now} and \texttt{.year}. After these cleaning and transformation steps are completed, each function returns a cleaned dataframe that is ready for further analysis. Finally, the cleaned dataframes are exported into individual CSV files for testing purposes, allowing the results of the data curation process to be reviewed and verified before being used in subsequent analysis.

Missing data values are handled during the data cleaning process to ensure that all dataframes contain consistent and usable values for analysis. When financial data is extracted from Yahoo Finance, missing values are typically represented by a double dash \texttt{(--)}. Since this placeholder is stored as a string, it can cause issues when performing numerical calculations or analysis on columns containing multiple data types. To address this, each cleaning function replaces instances of the double dash with an integer value of zero. This ensures that missing numerical values are represented consistently and that the columns can be converted into the appropriate data types for analysis. Additionally, any columns that contain no values are removed from the dataframe entirely, as they do not provide useful information for the analysis. These steps help maintain data consistency, reduce potential errors during processing, and ensure the resulting dataframes are suitable for further financial analysis.

The total number of attributes and records will vary depending on the stock and how much data is available for it at run-time, but averaging around 34 attributes and 160 records. A Discounted Cash Flow (DCF) model will be used for analysis of the stock and a Discounted Dividend Model (DDM) will be used for predictive analysis of future dividend payouts. These models use net income common stockholders, total assets, total liabilities, operating cash flow, capital expenditures, issuance of debt, repayment of debt, and dividend to calculate the equity value of a company's stock price. The attributes used for these models are in the dividend, income statement, cash flow, and balance sheet tables, and are listed below in the definitions table. Unused attributes may be useful in the future.

\begin{table}[ht]
\centering
\caption{Definitions of Financial Metrics}\label{tab:financial_metrics}
\begin{tabular}{p{4cm} p{8cm}}
\toprule
\textbf{Financial Metric} & \textbf{Definition} \\
\midrule
Net income common stockholders & The profit remaining for ordinary shareholders after deducting all operating expenses, taxes, and preferred stock dividends from a company's total net income. \\

Total assets & All resources owned or controlled by a company that have economic value. \\

Total liabilities & All financial debts and obligations owed to others. \\

Operating cash flow & The amount of cash generated by a company's core business operations over a specific period of time, excluding investing and financing activities. \\

Capital expenditures & The funds a company spends to acquire, upgrade, and maintain physical long-term assets. \\

Issuance of debt & The amount borrowed from investors in capital markets. \\

Repayment of debt & The amount paid toward outstanding debt owed to investors. \\

Dividends & A portion of a company's profits paid to shareholders. \\
\bottomrule
\end{tabular}
\end{table}

The Discounted Cash Flow (DCF) and Dividend Discount Model (DDM) are used to calculate the equity value of a stock based on a set of financial attributes. In this context, equity value serves as the dependent variable, as its calculated value is influenced by several independent variables, including net income common stockholders, total assets, total liabilities, operating cash flow, capital expenditures, debt issuance, debt repayment, and dividends. These financial attributes serve as independent variables because they represent individual components of a company's financial performance and position that contribute to the overall equity valuation. Although the values of these attributes may vary between companies and stocks, their combined influence determines the resulting equity value of each individual stock. By estimating the future equity value over a specified period, users can compare the projected value against the stock's current market price to help assess its potential financial risk and determine whether it aligns with their investment objectives.

*Note: move these to another, more relevent section later.

\begin{equation}
\begin{aligned}
&\textbf{Discounted Cash Flow (DCF) Model} \\[0.5em]
\text{Equity Value} &= \text{Present Value of Future Free Cash Flow to Equity} \\[1em]
    V_0 &= \frac{FCF_1^e}{(1+r_e)^1}
        + \frac{FCF_2^e}{(1+r_e)^2}
        + \frac{FCF_3^e}{(1+r_e)^3}
        + \cdots \\[1em]
    &= \sum_{t=1}^{\infty} \frac{FCF_t^e}{(1+r_e)^t}
\end{aligned}
\end{equation}

\begin{equation}
\begin{aligned}
&\textbf{Dividend Discount Model (DDM)} \\[0.5em]
\text{Equity Value} &= \text{Present Value of Expected Future Dividends} \\[1em]
V_0 &= \frac{div_1}{(1+r_e)^1}
     + \frac{div_2}{(1+r_e)^2}
     + \frac{div_3}{(1+r_e)^3}
     + \cdots \\[1em]
&= \sum_{t=1}^{\infty} \frac{div_t}{(1+r_e)^t}
\end{aligned}
\end{equation}

\section{Results}

Summary of analysis and data validation.

\section{Limitations}

Identify system issues and obstacles.
Possible solutions to overcome limitations.
Possibly any development obstacles that may have been encountered.

\section{Conclusion}

Reiterate goals/purpose of the project. Summarize results and findings.
Discuss future work and possible improvements.

\bibliographystyle{splncs04}
\bibliography{mybibliography}

% \subsection{Goals of this Research} 
% Specify exactly your aims of this paper. Also, write a sentence how you will address for each individual goal.

% \subsection{This is another subsection}
% Write something related to your paper.

% \section{Mobile Usability}


% \begin{table}
% \caption{Gas Prices}\label{gasprice}
% \begin{tabular}{|l|r|}
% \hline
% Month &  Price \\
% \hline
% Jan &  \$2.12 \\
% Feb &  0.12 \\
% Mar &  1.12 \\
% Apr &  2.12 \\
% May &  3.12 \\
% Jun &  1.12 \\

% \hline
% \end{tabular}
% \end{table}



% \begin{table}
% \caption{Table captions should be placed above the
% tables.}\label{tab1}
% \begin{tabular}{|l|l|l|}
% \hline
% Heading level &  Example & Font size and style\\
% \hline
% Title (centered) &  {\Large\bfseries Lecture Notes} & 14 point, bold\\
% 1st-level heading &  {\large\bfseries 1 Introduction} & 12 point, bold\\
% 2nd-level heading & {\bfseries 2.1 Printing Area} & 10 point, bold\\
% 3rd-level heading & {\bfseries Run-in Heading in Bold.} Text follows & 10 point, bold\\
% 4th-level heading & {\itshape Lowest Level Heading.} Text follows & 10 point, italic\\
% \hline
% \end{tabular}
% \end{table}


% The gasoline prices are summarized in Table \ref{gasprice}



% \noindent Displayed equations are centered and set on a separate
% line.
% \begin{equation}
% x + y = z
% \end{equation}
% Please try to avoid rasterized images for line-art diagrams and
% schemas. Whenever possible, use vector graphics instead (see
% Fig.~\ref{fig1}).

% \begin{figure}
% \includegraphics[width=\textwidth]{fig1.eps}
% \caption{A figure caption is always placed below the illustration.
% Please note that short captions are centered, while long ones are
% justified by the macro package automatically.} 
% \label{fig1}
% \end{figure}

% The detailed steps of the MEC is shown in Figure \ref{mecFig}.

% \begin{figure}
% \includegraphics[width=\textwidth]{mecomputing.png}
% \caption{This is the MEC} 
% \label{mecFig}
% \end{figure}

% %\begin{theorem}
% %This is a sample theorem. The run-in heading is set in bold, while
% %the following text appears in italics. Definitions, lemmas,
% %propositions, and corollaries are styled the same way.
% %\end{theorem}
% %
% % the environments 'definition', 'lemma', 'proposition', 'corollary',
% % 'remark', and 'example' are defined in the LLNCS documentclass as well.
% %
% \begin{proof}
% Proofs, examples, and remarks have the initial word in italics,
% while the following text appears in normal font.
% \end{proof}


% The following are the items in my lunch menu.
% \begin{itemize}
%     \item Fish
%     \begin{itemize}
%         \item sauce
%         \item fires
%     \end{itemize}
%     \item Chicken
%     \item Yogurt
%     \item ice cream
% \end{itemize}

% The following are the steps to cook this recipe. 



% \begin{enumerate}
% \item Clean vessels
% \item Cook
% \begin{enumerate}
%     \item On the stove
%     \item Place the pan
%     \item Crack the egg
% \end{enumerate}
% \item Eat
% \end{enumerate}



% For citations of references, we prefer the use of square brackets
% and consecutive numbers. \cite{bariah2020prospective}, \cite{kato2020ten}.

% Now I am citing this article. \cite{hoehle2015mobile}

% citing now \cite{8766917}, \cite{10.5555/2667432.2667458}

% Now, I am using my sixth \cite{mao2017survey} citation.

% \cite{910896}

% \cite{*}

%
% ---- Bibliography ----
%
% BibTeX users should specify bibliography style 'splncs04'.
% References will then be sorted and formatted in the correct style.

\end{document}

